% Options for packages loaded elsewhere
\PassOptionsToPackage{unicode}{hyperref}
\PassOptionsToPackage{hyphens}{url}
\PassOptionsToPackage{dvipsnames,svgnames,x11names}{xcolor}
%
\documentclass[
]{article}

\usepackage{amsmath,amssymb}
\usepackage{iftex}
\ifPDFTeX
  \usepackage[T1]{fontenc}
  \usepackage[utf8]{inputenc}
  \usepackage{textcomp} % provide euro and other symbols
\else % if luatex or xetex
  \usepackage{unicode-math}
  \defaultfontfeatures{Scale=MatchLowercase}
  \defaultfontfeatures[\rmfamily]{Ligatures=TeX,Scale=1}
\fi
\usepackage{lmodern}
\ifPDFTeX\else  
    % xetex/luatex font selection
\fi
% Use upquote if available, for straight quotes in verbatim environments
\IfFileExists{upquote.sty}{\usepackage{upquote}}{}
\IfFileExists{microtype.sty}{% use microtype if available
  \usepackage[]{microtype}
  \UseMicrotypeSet[protrusion]{basicmath} % disable protrusion for tt fonts
}{}
\makeatletter
\@ifundefined{KOMAClassName}{% if non-KOMA class
  \IfFileExists{parskip.sty}{%
    \usepackage{parskip}
  }{% else
    \setlength{\parindent}{0pt}
    \setlength{\parskip}{6pt plus 2pt minus 1pt}}
}{% if KOMA class
  \KOMAoptions{parskip=half}}
\makeatother
\usepackage{xcolor}
\setlength{\emergencystretch}{3em} % prevent overfull lines
\setcounter{secnumdepth}{5}
% Make \paragraph and \subparagraph free-standing
\makeatletter
\ifx\paragraph\undefined\else
  \let\oldparagraph\paragraph
  \renewcommand{\paragraph}{
    \@ifstar
      \xxxParagraphStar
      \xxxParagraphNoStar
  }
  \newcommand{\xxxParagraphStar}[1]{\oldparagraph*{#1}\mbox{}}
  \newcommand{\xxxParagraphNoStar}[1]{\oldparagraph{#1}\mbox{}}
\fi
\ifx\subparagraph\undefined\else
  \let\oldsubparagraph\subparagraph
  \renewcommand{\subparagraph}{
    \@ifstar
      \xxxSubParagraphStar
      \xxxSubParagraphNoStar
  }
  \newcommand{\xxxSubParagraphStar}[1]{\oldsubparagraph*{#1}\mbox{}}
  \newcommand{\xxxSubParagraphNoStar}[1]{\oldsubparagraph{#1}\mbox{}}
\fi
\makeatother


\providecommand{\tightlist}{%
  \setlength{\itemsep}{0pt}\setlength{\parskip}{0pt}}\usepackage{longtable,booktabs,array}
\usepackage{calc} % for calculating minipage widths
% Correct order of tables after \paragraph or \subparagraph
\usepackage{etoolbox}
\makeatletter
\patchcmd\longtable{\par}{\if@noskipsec\mbox{}\fi\par}{}{}
\makeatother
% Allow footnotes in longtable head/foot
\IfFileExists{footnotehyper.sty}{\usepackage{footnotehyper}}{\usepackage{footnote}}
\makesavenoteenv{longtable}
\usepackage{graphicx}
\makeatletter
\def\maxwidth{\ifdim\Gin@nat@width>\linewidth\linewidth\else\Gin@nat@width\fi}
\def\maxheight{\ifdim\Gin@nat@height>\textheight\textheight\else\Gin@nat@height\fi}
\makeatother
% Scale images if necessary, so that they will not overflow the page
% margins by default, and it is still possible to overwrite the defaults
% using explicit options in \includegraphics[width, height, ...]{}
\setkeys{Gin}{width=\maxwidth,height=\maxheight,keepaspectratio}
% Set default figure placement to htbp
\makeatletter
\def\fps@figure{htbp}
\makeatother

\usepackage{booktabs}
\usepackage{longtable}
\usepackage{array}
\usepackage{multirow}
\usepackage{wrapfig}
\usepackage{float}
\usepackage{colortbl}
\usepackage{pdflscape}
\usepackage{tabu}
\usepackage{threeparttable}
\usepackage{threeparttablex}
\usepackage[normalem]{ulem}
\usepackage{makecell}
\usepackage{xcolor}
\usepackage{booktabs}
\usepackage{longtable}
\usepackage[a4paper, left=0.75in, right=1in, top=1in, bottom=1in]{geometry}
\makeatletter
\@ifpackageloaded{caption}{}{\usepackage{caption}}
\AtBeginDocument{%
\ifdefined\contentsname
  \renewcommand*\contentsname{Table of contents}
\else
  \newcommand\contentsname{Table of contents}
\fi
\ifdefined\listfigurename
  \renewcommand*\listfigurename{List of Figures}
\else
  \newcommand\listfigurename{List of Figures}
\fi
\ifdefined\listtablename
  \renewcommand*\listtablename{List of Tables}
\else
  \newcommand\listtablename{List of Tables}
\fi
\ifdefined\figurename
  \renewcommand*\figurename{Figure}
\else
  \newcommand\figurename{Figure}
\fi
\ifdefined\tablename
  \renewcommand*\tablename{Table}
\else
  \newcommand\tablename{Table}
\fi
}
\@ifpackageloaded{float}{}{\usepackage{float}}
\floatstyle{ruled}
\@ifundefined{c@chapter}{\newfloat{codelisting}{h}{lop}}{\newfloat{codelisting}{h}{lop}[chapter]}
\floatname{codelisting}{Listing}
\newcommand*\listoflistings{\listof{codelisting}{List of Listings}}
\makeatother
\makeatletter
\makeatother
\makeatletter
\@ifpackageloaded{caption}{}{\usepackage{caption}}
\@ifpackageloaded{subcaption}{}{\usepackage{subcaption}}
\makeatother

\ifLuaTeX
  \usepackage{selnolig}  % disable illegal ligatures
\fi
\usepackage{bookmark}

\IfFileExists{xurl.sty}{\usepackage{xurl}}{} % add URL line breaks if available
\urlstyle{same} % disable monospaced font for URLs
\hypersetup{
  pdftitle={Differential abundance analysis},
  colorlinks=true,
  linkcolor={blue},
  filecolor={Maroon},
  citecolor={Blue},
  urlcolor={Blue},
  pdfcreator={LaTeX via pandoc}}


\title{Differential abundance analysis}
\author{}
\date{}

\begin{document}
\maketitle


\subsection{Differential abundance
analysis}\label{differential-abundance-analysis}

Analysis information:

\begin{itemize}
\tightlist
\item
  Taxonomic level: species
\end{itemize}

Taxa that are less prevalent than this (in relabundance) have been
agglomerated:

\begin{itemize}
\tightlist
\item
  Detection threshold: 0.1\%
\item
  Prevalence threshold: 0.1
\item
  Number of taxonomic groups: 252
\end{itemize}

BMI is not significantly different between conditions or time points.
Hence we can ignore it in the differential abundance testing. In group
``B'' it seems to be borderline significant (p\textless0.05 before
multiple testing adjustment).

Number of significant findings within time points.

\begin{itemize}
\tightlist
\item
  Kruskal test, Benjamini-Hochberg FDR adjustment
\end{itemize}

\begin{verbatim}
before  after 
     1      8 
\end{verbatim}

Number of significant findings within treatments:

\begin{itemize}
\tightlist
\item
  Paired Wilcoxon test, Benjamini-Hochberg FDR adjustment
\end{itemize}

\begin{verbatim}
rice  oat 
   9    7 
\end{verbatim}

Let us look more closely to the significant hits. The logFC corresponds
to logFC in CLR (tp2-tp1).

\begin{figure}

\begin{minipage}{0.33\linewidth}
\includegraphics{daa_species_files/figure-pdf/hits-1.pdf}\end{minipage}%
%
\begin{minipage}{0.33\linewidth}
\includegraphics{daa_species_files/figure-pdf/hits-2.pdf}\end{minipage}%
%
\begin{minipage}{0.33\linewidth}
\includegraphics{daa_species_files/figure-pdf/hits-3.pdf}\end{minipage}%
\newline
\begin{minipage}{0.33\linewidth}
\includegraphics{daa_species_files/figure-pdf/hits-4.pdf}\end{minipage}%
%
\begin{minipage}{0.33\linewidth}
\includegraphics{daa_species_files/figure-pdf/hits-5.pdf}\end{minipage}%
%
\begin{minipage}{0.33\linewidth}
\includegraphics{daa_species_files/figure-pdf/hits-6.pdf}\end{minipage}%
\newline
\begin{minipage}{0.33\linewidth}
\includegraphics{daa_species_files/figure-pdf/hits-7.pdf}\end{minipage}%
%
\begin{minipage}{0.33\linewidth}
\includegraphics{daa_species_files/figure-pdf/hits-8.pdf}\end{minipage}%
%
\begin{minipage}{0.33\linewidth}
\includegraphics{daa_species_files/figure-pdf/hits-9.pdf}\end{minipage}%
\newline
\begin{minipage}{0.33\linewidth}
\includegraphics{daa_species_files/figure-pdf/hits-10.pdf}\end{minipage}%
%
\begin{minipage}{0.33\linewidth}
\includegraphics{daa_species_files/figure-pdf/hits-11.pdf}\end{minipage}%
%
\begin{minipage}{0.33\linewidth}
\includegraphics{daa_species_files/figure-pdf/hits-12.pdf}\end{minipage}%
\newline
\begin{minipage}{0.33\linewidth}
\includegraphics{daa_species_files/figure-pdf/hits-13.pdf}\end{minipage}%
%
\begin{minipage}{0.33\linewidth}
\includegraphics{daa_species_files/figure-pdf/hits-14.pdf}\end{minipage}%
%
\begin{minipage}{0.33\linewidth}
\includegraphics{daa_species_files/figure-pdf/hits-15.pdf}\end{minipage}%
\newline
\begin{minipage}{0.33\linewidth}
\includegraphics{daa_species_files/figure-pdf/hits-16.pdf}\end{minipage}%

\end{figure}%

\begin{longtable}[]{@{}
  >{\raggedright\arraybackslash}p{(\columnwidth - 10\tabcolsep) * \real{0.0612}}
  >{\raggedright\arraybackslash}p{(\columnwidth - 10\tabcolsep) * \real{0.3776}}
  >{\raggedleft\arraybackslash}p{(\columnwidth - 10\tabcolsep) * \real{0.1020}}
  >{\raggedleft\arraybackslash}p{(\columnwidth - 10\tabcolsep) * \real{0.0714}}
  >{\raggedleft\arraybackslash}p{(\columnwidth - 10\tabcolsep) * \real{0.1939}}
  >{\raggedleft\arraybackslash}p{(\columnwidth - 10\tabcolsep) * \real{0.1939}}@{}}
\toprule\noalign{}
\begin{minipage}[b]{\linewidth}\raggedright
Group
\end{minipage} & \begin{minipage}[b]{\linewidth}\raggedright
Taxon
\end{minipage} & \begin{minipage}[b]{\linewidth}\raggedleft
logFC.clr
\end{minipage} & \begin{minipage}[b]{\linewidth}\raggedleft
FDR
\end{minipage} & \begin{minipage}[b]{\linewidth}\raggedleft
mean.relabundance1
\end{minipage} & \begin{minipage}[b]{\linewidth}\raggedleft
mean.relabundance2
\end{minipage} \\
\midrule\noalign{}
\endhead
\bottomrule\noalign{}
\endlastfoot
rice & Alistipes\_communis & -0.0016 & 0.2430 & 0.295 & 0.135 \\
rice & Bifidobacterium\_longum & -0.0049 & 0.0020 & 1.794 & 1.311 \\
rice & Bilophila\_wadsworthia & 0.0003 & 0.1263 & 0.066 & 0.096 \\
rice & Candidatus\_Cibionibacter\_quicibialis & -0.0147 & 0.0003 & 2.336
& 0.822 \\
rice & Clostridiaceae\_bacterium & 0.0075 & 0.0277 & 0.648 & 1.422 \\
rice & Eubacterium\_rectale & -0.0127 & 0.1450 & 3.214 & 1.878 \\
rice & Eubacterium\_siraeum & -0.0006 & 0.2333 & 0.262 & 0.205 \\
rice & Firmicutes\_bacterium\_AF16\_15 & 0.0031 & 0.0790 & 0.317 &
0.627 \\
rice & GGB79734\_SGB15291 & 0.0006 & 0.1263 & 0.116 & 0.178 \\
oat & Blautia\_massiliensis & 0.0097 & 0.1564 & 1.169 & 2.183 \\
oat & Candidatus\_Cibionibacter\_quicibialis & -0.0052 & 0.1564 & 1.548
& 1.014 \\
oat & Clostridiales\_bacterium\_KLE1615 & -0.0025 & 0.1711 & 0.667 &
0.412 \\
oat & GGB27106\_SGB6188 & 0.0010 & 0.1753 & 0.019 & 0.115 \\
oat & GGB4569\_SGB6310 & 0.0004 & 0.1564 & 0.029 & 0.070 \\
oat & GGB9559\_SGB14969 & 0.0002 & 0.1564 & 0.031 & 0.049 \\
oat & Lawsonibacter\_asaccharolyticus & -0.0005 & 0.1753 & 0.164 &
0.111 \\
\end{longtable}




\end{document}
